\documentclass[10pt]{article}
\usepackage[a4paper,margin=20mm]{geometry}
\usepackage{amsmath,amssymb,amsthm}
\usepackage[T1]{fontenc}
\usepackage{hyperref}
\newtheorem{theorem}{Theorem}
\title{Conjecture 00000001255 is false\\A noisy one-site cellular automaton}
\author{GodBlf}
\date{October 8, 2026}
\begin{document}
\maketitle
\begin{abstract}
A binary flip rule on a one-site periodic ring has no deterministic fixed
configuration. Its nondegenerate noisy update has a unique stationary law,
the uniform law. This law is an extreme stationary measure and is not a
point mass. Thus the conjecture's extreme-point assertion fails.
\end{abstract}
\section{Statement and model}
The conjecture asserts, for a cellular automaton on a finite ring driven by
Glauber-type noise, that the relevant stationary simplex has deterministic
fixed points as its extreme points. It also asserts a support description
and a dimension identity. A counterexample to the extreme-point clause
suffices to refute their conjunction.

Take the periodic ring $\mathbb Z/\mathbb Z$, with binary state space
$S=\{0,1\}$. The local deterministic rule is $F(s)=1-s$.
At an update, independently flip the prescribed output with probability
$1/4$. Equivalently, apply $F$ with probability $3/4$, and keep the current
bit with probability $1/4$. This is a one-site, positive-temperature
Glauber-type noisy local update. In the order $(0,1)$ its transition matrix is
\[
 P=\begin{pmatrix}1/4&3/4\\3/4&1/4\end{pmatrix}.
\]
Every entry is positive and each row sums to one. The deterministic rule
has a two-cycle and no fixed point.

\section{Exact stationary and extreme-point computation}
\begin{theorem}
The only stationary probability law is $u=(1/2,1/2)$. It is an extreme
point of the convex set of stationary laws, and is not a Dirac law.
\end{theorem}
\begin{proof}
Write a probability law as $(a,b)$, with $a,b\geq0$ and $a+b=1$.
Stationarity at state zero gives
$a=a/4+3b/4$, hence $a=b=1/2$. Conversely, direct multiplication gives
$uP=u$. The stationary-law set is therefore the singleton $\{u\}$.
If $u=r v+(1-r)w$ with stationary $v,w$ and $0<r<1$, uniqueness implies
$v=w=u$, which is exactly the extreme-point condition. Both coordinates
of $u$ are $1/2$, so it is neither of the two Dirac laws.
\end{proof}

The unique extreme stationary law cannot be supported at a deterministic
fixed point because there are none. If the statement instead intends the
simplex of probability laws on the stationary support, that support is all
of $S$; its vertices are the two Dirac laws, and neither underlying
configuration is a fixed point of $F$. Thus that reading fails as well.
No interpretation of the undefined ``minimal attractor dimension'' is
needed for this disproof.

\section{Formal verification}
The accompanying Lean project represents the two configurations by
\texttt{Bool}, and probability laws by functions $S\to\mathbb R$ with
nonnegative coordinates summing to one. This represents every probability
measure on the finite state space, not a restricted numerical sample.
The positivity and row-sum theorems verify the
kernel. \texttt{stationary\_unique} quantifies over all real-valued laws;
\texttt{uniform\_extreme} uses the standard convex-combination definition.
\texttt{conjecture\_false} negates the extreme-point clause, using the actual
flip rule and the actual stationary set.

Lean and Mathlib are pinned by the toolchain and Lake manifest.
Run \texttt{lake build} and \texttt{lake env lean Check.lean} in the
\texttt{lean} directory. No admitted proofs, additional axioms, or
\texttt{native\_decide} are used. The axiom audit is in
\texttt{verification.txt}. Rebuild this document with
\texttt{pdflatex main.tex} twice.
\end{document}
